\section{Introduction}\label{sec:intro}

% Pinned to the top of page 2: declared here, `[t]`-only, and the ACL title block
% owns the top of page 1, so this is the first top slot LaTeX can use. Do not add
% `b`/`h` --- they would let it drift.
\begin{figure}[t]
\centering
\includegraphics[width=0.90\textwidth]{fig_overview}
\caption{What this paper measures. \textbf{(1)} One model call per step emits a thought, a grounding and an action, with the last eight action outcomes as context. \textbf{(2)} That step is grounded by element id, by numbered mark, or by coordinate; screenshot-free variants separate the text payload from the prompt family (\S\ref{sec:setup}). \textbf{(3)} Three of our five routing policies, bracketed by an outcome oracle above and a fixed single mode below (the pooled-tier router and the zero-token rule are omitted for space; \S\ref{sec:lowerbound}), and the two ceilings a perfect per-task choice could reach. Only the cost ceiling survives the rerun control of \S\ref{sec:noise}, because it adds no arm. \textbf{Caution:} the policy marked as exceeding a fixed single mode does so on a lossless-saving criterion under a label-shuffle null, and the one setting surviving correction is \texttt{red$\cdot$B2}, whose triage signal is below chance (AUROC 0.483; Appendix Table~\ref{tab:pb05}). Under the stricter test no cell Pareto-dominates always-cheapest (Appendix Table~\ref{tab:pb06}).}
\label{fig:overview}
\end{figure}

A web agent has to be shown the page before it can act on it. In practice it is shown one of three things: the accessibility tree as text, a screenshot as pixels, or a set-of-marks fusion in which the screenshot is annotated with the tree's element ids \citep{yang2023som}. Which of these an agent receives is the grounding decision that \citet{zheng2024seeact} identify as the bottleneck for a capable backbone, and that a line of work on visual grounding for GUI agents attacks directly \citep{zheng2024uground}. Benchmarks such as VisualWebArena \citep{koh2024visualwebarena}, WebArena \citep{zhou2024webarena} and Mind2Web \citep{deng2023mind2web} report agents that make this choice once, at build time, and hold it for every task. The choice is treated as a design decision about the system rather than as a decision about the task in front of it.

\looseness=-1 It need not be. We hold one scaffold, one prompt budget and one action space fixed, and compare the six representations across eight (site $\times$ backbone) cells. They are not redundant: 44 of the 48 mode--cell pairs solve at least one task no other mode in that cell solved, and every mode does so in some cell. Where exactly one channel succeeds, the losing channel fails in a structurally different way. \emph{Which} representation wins also reverses between our two task sets, so the winner is a property of the deployment and not of the modality (\S\ref{sec:complementarity}). A per-task choice therefore has something to choose between, which raises the obvious next question: what would it be worth?

\looseness=-1 That question is easy to answer badly. An outcome oracle over six modes (it sees every outcome and picks a winning mode per task) reaches 7.1--51.9\% against a best single mode of 2.2--35.6\%, which reads as enormous. It is also a six-arm union quoted against one arm: a task counts as solved there if \emph{any} of six runs solved it, and an \emph{arm} is one mode run in one cell. The comparison it needs is a rerun. Rerunning one condition changes 12--14\% of task outcomes, and 22\% of tasks flip in at least one of three replicated arms (both measured in two cells and assumed to transfer to the other six; Appendix~\ref{app:limitations}); at that rate a second copy of an arm already in hand raises the union ceiling by as much as a different representation does. The union survives as a measured bound; what does not is the attribution of its headroom to representation diversity. What holds up without adding an arm is narrower: routing only the never-solved tasks to the cheapest mode cuts cost 9.5--30.6\% in every cell at unchanged success (\S\ref{sec:upperbound}).

Routing itself is not new. Systems route a query to one of several \emph{models} to buy accuracy at lower cost \citep{chen2023frugalgpt,ong2025routellm}, including for web and computer-use agents specifically \citep{webrouter2025,liu2026adaptive}; \citet{moslem2026routingsurvey} survey the area. What those route over is a capability gradient: a larger model really is more accurate, so difficulty predicts which arm to use. We ask whether the same holds when the arms are representations of one browser state under one model, where no such gradient exists by construction.

The lower bound is where the paper turns. We build five routing policies: choosing which mode to use, a learned triage of spend, a zero-token rule read off the task text, a confidence cascade, and pooled cost tiers. None yields a robust improvement over simply fixing the best or cheapest arm, the one exception being descriptive rather than confirmatory (\S\ref{sec:lowerbound}). That would be uninformative if they all failed the same way, and they do not. One dies of label supply. One has ample labels and buys nothing. One has a perfect free signal and still loses, because the expensive arm does not hurt on the tasks the signal excludes (\S\ref{sec:lowerbound}).

The gap between those bounds is what this paper can offer without overstating it. What it comes down to is a circularity we call the \textbf{supply--value coupling}: routing supervision is produced at the success rate, so routing is least learnable exactly where it would be most valuable. Those look like two obstructions and are one: the labels a router can learn from and the tasks where routing can pay both track the agent's success rate, which indexes the negative result to the current capability regime and says what would overturn it (\S\ref{sec:gap}).

Underneath all of it, every comparison is read against a measured same-condition rerun band rather than against zero (\S\ref{sec:noise}); \citet{hajimiri2026budgetmatched} make the same move on a different axis, showing that a token-matched vanilla baseline erases the reported gains of several agent memory and skill modules. Every quantity is read from its producing artefact at render time (Appendix~\ref{sec:threats}). On workloads at these success rates the resolution of the instrument is comparable to the effects being reported, and a paper that does not say so is reporting its instrument.

\looseness=-1 \paragraph{Reading this as a statement about visual grounding.} Four of the six modes carry no screenshot, and two of those four replace the accessibility tree with the flattened marks that annotate one (\S\ref{sec:setup}). Only one pair isolates the pixels while holding text payload, prompt family and element-id regime fixed; every other image-on/image-off contrast moves the identifier regime with it, so this design bounds what the visual channel is worth rather than cleanly ablating it. What the wider set does establish is that \emph{which} channel wins is a property of the deployment and not of the modality (\S\ref{sec:complementarity}). The part that transfers is the null: on workloads at these success rates, a grounding intervention worth fewer points than a same-condition rerun of the arm it modifies is not yet distinguishable from one (\S\ref{sec:noise}).

\paragraph{Contributions.} (i) A \emph{rerun-matched} way to price any multi-arm ceiling: the null for a union or oracle ceiling is a same-condition replicate, not zero, and we derive the resolution threshold that follows (\S\ref{sec:noise}). (ii) Applied here, the six representations are genuinely complementary and which one wins is a property of the deployment (\S\ref{sec:complementarity}), but most of the apparent oracle headroom is rerun variance; what survives the correction is the cost ceiling, 9.5--30.6\% at unchanged success (\S\ref{sec:upperbound}). (iii) Five routing constructions all land at or below trivial fixed policies, and the obstruction is the supervision rather than the estimator: labels arrive at the agent's success rate (\S\ref{sec:lowerbound}--\S\ref{sec:gap}).
